% !TeX root = main.tex
\lecture{6}{Fri 30 Jan 2026 12:00}{Partial Differentiation VI: 3D Stationary Points}

\section{2D Recap}
In 2D, a stationary point of a function $f(x)$ occurs at $x = a$ where $f(x)$ is neither increasing or decreasing such that $f^\prime(a) = 0$.

Consider the Taylor series about $x = a$, where $f^\prime(a) = 0$:

\[
    f(a+h) = f(a) + h f^\prime(a) + \frac{h^2}{2!} f^{\prime\prime} (a) + \cdots = f(a) + \frac{h^2}{2!} f^{\prime\prime} (a) + \cdots
\]

If $f^{\prime\prime} (a) > 0$, the function looks like $+h^2$ near the stationary point. This gives us a minima. If the second derivative is negative the function looks like $-h^2$, giving a maximum point.

If the second derivative is zero, we need to find the first non-zero derivative $f^n{(a)} \neq 0$ where $n > 2$.
\begin{itemize}
    \item If $n = 2k + 1$ (odd), the function looks like $\pm h^{2k + 1}$ near the stationary point. This results in a point of inflexion.
    \item If $n = 2k$ (even) and $f^{2k}(a) > 0$, the function looks like $+ h^{2k}$ near the stationary point so that we have a minimum.
    \item If $n = 2k$ (even) and $f^{2k}(a) < 0$, the function looks like $- h^{2k}$ near the stationary point so that we have a maximum.
\end{itemize}

\section{Stationary Points of Several Variables}
\subsection{Two Variables}
Consider a function $f(x, y)$ at a point $(a, b)$. For a stationary point to occur here, the function must be flat in all directions at that point, so all the partial derivatives must be zero:
\[
    f_x(a, b) = f_y(a, b) = 0
\]

To classify it, we again consider the Taylor series:
\[
    f(a+h, b+k) \approx f(a, b) + \frac{1}{2} \left(Ah^2 + 2Bhk + Ck^2\right)
\]
Where $A = f_{xx}(a, b)$, $B = f_{xy}(a,b)$ and $C = f_{yy}(a,b)$. We can then complete the square:

\[
    Ah^2 + Bhk + Ck^2 = A \left(h + \frac{B}{A}k\right)^2 + \left(C - \frac{B^2}{A}\right)k^2
\]

\begin{itemize}
    \item For the stationary point to be a minimum, we need $A > 0$ and $C - B^2 / A = AC - B^2 > 0$.
    \item For a maximum, we need $A < 0$ and $C-B^2/A = AC-B^2 < 0$.
    \item For a saddle point, with a minimum in one direction and a maximum in the other, we need $A$ and $C - B^2/A$ to be oppositely signed. This corresponds to the case where $AC - B^2 > 0$.
\end{itemize}

Considering the second derivatives directly:
\begin{itemize}
    \item For a minimum we need $f_{xx} > 0$ and the Hessian determinant must be positive:
    \[
        \Delta = \det H = \det \begin{bmatrix}
            f_{xx} & f_{xy} \\
            f_{yx} & f_{yy}
        \end{bmatrix} > 0
    \]

    \item For a maximum, we need $f_{xx} < 0$ and the Hessian determinant must be positive:
    \[
        \Delta = \det H = \det \begin{bmatrix}
            f_{xx} & f_{xy} \\
            f_{yx} & f_{yy}
        \end{bmatrix} > 0
    \]
    
    \item For a saddle point, we need the Hessian determinant to be negative:
    \[
        \Delta = \det H = \det \begin{bmatrix}
            f_{xx} & f_{xy} \\
            f_{yx} & f_{yy}
        \end{bmatrix} < 0
    \] 
\end{itemize}

\subsection{Generalising To 3D}
For a function $f(x, y, z)$ to have a stationary point at point $(a, b, c)$, it must be flat in all directions around that point. Therefore, all the first derivatives must be zero:

\[
    f_x(a, b, c) = f_y(a, b, c) = f_z(a, b, c) = 0
\]

The Taylor Series about the point $(a, b, c)$ is:
\[
    f(a+h, b+k, c+l) \approx f(a, b, c) + \frac{1}{2} \left(Ah^2 + 2Bhk + 2Chl+Dk^2 + 2Ekl + Fl^2\right)
\] 

Where $A = f_{xx}, B = f_{xy}, C = f_{xz}, D = f_{yy}, E = f_{yz}, F = f_{zz}$. Again we complete the square:

\[
    A \left(h + \frac{B}{A}k + \frac{C}{A}l\right)^2 + \left(D - \frac{B^2}{A}\right)k^2 + 2 \left(E - \frac{BC}{A}\right)kl + \left(F - \frac{C^2}{A}\right)l^2
\]
